\chapter{Введение}\label{ch:intro}

% TODO (M2): Развёрнутое введение.
% Структура (по Chapter Plan, RESEARCH.md → § 1, ~3 страницы):
%   1.1 Актуальность темы
%       — массовое применение EPD (e-readers, ценники, signage);
%       — конфликт «энергия ↔ ghosting ↔ скорость ↔ ресурс»;
%       — состояние области (Kang 2025, Lai 2026 — современные алгоритмы);
%       — пробел: нет работ со структурной теоремой об оптимальной форме waveform.
%   1.2 Цель и задачи работы
%   1.3 Научная новизна
%       — N2+N1 гибрид: алгоритм + структурная теорема (bang-bang с ≤K_eff фаз);
%       — каскадная теория PDE → ODE → дискретный PMP → Парето;
%       — экспериментальная валидация на стенде.
%   1.4 Положения, выносимые на защиту (3-4 пункта)
%   1.5 Апробация (если будет — конференция / препринт)
%   1.6 Структура работы

\section*{Заглушка}

Раздел будет наполнен на этапе~M2 в соответствии с Chapter Plan
(см. \texttt{.ai-factory/RESEARCH.md}).

Тест цитирования (будет удалён на~M2): современные обзоры предметной
области~\cite{zhong2026_review,heikenfeld2011_review}, базовая физика
электрофореза~\cite{comiskey1998_nature,yang2021_mechanisms}, прямые
конкуренты по waveform-оптимизации~\cite{kang2025_dp_waveform,lai2026_e3dw_energy}.
Используемая математическая база~\cite{pontryagin1983_optimal_processes,vasiliev2002_methods,phogat2017_discrete_pmp}.
